\section{Arquitetura Docker}

\begin{frame}
  \frametitle{Arquitetura Cliente-Servidor}
  \begin{block}{Visão Geral}
    O \textit{Docker client} comunica-se com o \textit{Docker daemon}, que executa o \textit{build}, execução e distribuição dos containers.
  \end{block}

  \vspace{0.5em}
  \begin{itemize}
    \item<2-> Daemon pode estar no mesmo sistema ou em um sistema remoto
    \item<3-> Comunicação via REST API, usando UNIX \textit{sockets} ou uma interface de rede
    \item<4-> Docker Compose também é um cliente Docker, permite trabalhar com aplicações que utilizam vários containers
  \end{itemize}
\end{frame}


\begin{frame}
  \frametitle{Docker Daemon/dockerd}
  Escuta os \textit{requests} da Docker API e gerencia os objetos Docker.

  \vspace{0.5em}
  Objetos gerenciados:
  \begin{itemize}
    \item Imagens
    \item Containers
    \item \textit{Networks} e Volumes
  \end{itemize}

  \vspace{0.5em}
  Um daemon também pode se comunicar com outros daemons para gerenciar serviços.
\end{frame}


\begin{frame}
  \frametitle{Docker Client}
    Forma primária de interagir com o Docker, envia comandos ao \texttt{dockerd} via Docker API.

  \vspace{0.5em}
  \begin{itemize}
    \item Exemplo: \texttt{docker run} → cliente envia ao daemon → daemon processa
    \item Um cliente pode se comunicar com \textbf{mais de um daemon}
  \end{itemize}
\end{frame}


\begin{frame}
  \frametitle{Docker Desktop}
  Aplicação para \textbf{Mac, Windows e Linux} para buildar e compartilhar containers.

  \vspace{0.5em}
  Inclui:
  \begin{itemize}
    \item Docker daemon (\texttt{dockerd}) e Docker client (\texttt{docker})
    \item Docker Compose e Docker Content Trust
    \item Kubernetes e Credential Helper
  \end{itemize}
\end{frame}


\begin{frame}
  \frametitle{Docker Registries}
    Repositórios para armazenar e distribuir Docker images.

  \vspace{0.5em}
  \begin{itemize}
    \item \textbf{Docker Hub} é o registro público padrão
    \item Também é possível manter registros \textbf{privados}
  \end{itemize}

  \vspace{0.5em}
  \onslide<4->{
    \begin{columns}
      \column{0.33\textwidth}
        \texttt{docker pull} \\ obtém a imagem
      \column{0.33\textwidth}
        \texttt{docker run} \\ obtém e executa
      \column{0.33\textwidth}
        \texttt{docker push} \\ envia a imagem
    \end{columns}
  }
\end{frame}


\begin{frame}
  \frametitle{Docker Objects}
  Ao usar o Docker, trabalha-se com imagens, containers, \textit{networks},
  volumes, \textit{plugins} e outros objetos.
\end{frame}

\begin{frame}
  \frametitle{Docker Objects — \textit{Image}}
  \begin{block}{O que é uma \textit{image}?}
    \textit{Template} com instruções para criar um container. Geralmente baseada em
    outra image, com customizações adicionais.
  \end{block}

  \vspace{0.5em}
  \begin{itemize}
    \item Definida por um \textbf{Dockerfile}
    \item Cada instrução do Dockerfile cria uma \textbf{camada} na image
    \item Ao refazer o \textit{build}, só as camadas \textbf{modificadas} são reconstruídas
    \item Resultado: images leves, pequenas e rápidas
  \end{itemize}
\end{frame}

\begin{frame}
  \frametitle{Docker Objects — \textit{Containers}}
  \begin{block}{O que é um container?}
    Uma \textbf{instância em execução} de uma \textit{image}.
  \end{block}

  \vspace{0.5em}
  \begin{itemize}
    \item É possível criar, iniciar, parar, mover ou deletar via API ou CLI
    \item Conectável a \textit{networks} e volumes de armazenamento
    \item Isolado por padrão, nível de isolamento configurável
    \item \alert{Mudanças não persistidas são perdidas quando remove o container}
  \end{itemize}
\end{frame}

\begin{frame}
  \frametitle{Exemplo: \texttt{docker run}}
  \texttt{docker run -i -t ubuntu /bin/bash}

  \vspace{0.5em}
  O que acontece ao executar:

  \begin{enumerate}
    \item<2-> Baixa a \textit{image} \texttt{ubuntu} do registro, se não existir localmente
    \item<3-> Cria um novo container
    \item<4-> Aloca sistema de arquivos de leitura e escrita
    \item<5-> Cria interface de rede e atribui IP ao container
    \item<6-> Inicia o container e executa \texttt{/bin/bash} interativamente
    \item<7-> Ao executar \texttt{exit} — container \textbf{para}, mas \textbf{não é removido}
  \end{enumerate}
\end{frame}